\paragraph{完整流程比较。}\AchievementDetectionZH{} 提示分组配对不确定性见表~\ref{tab:uncertainty}。

\begin{table}[t]
\centering\footnotesize
\caption{域内检测结果（百分比）。报告可用种子的均值与样本标准差；单种子不显示标准差。阈值仅依据域内校准集确定。}\label{tab:main}
\begin{tabular}{lrrrr}\toprule
Method & R@128 & FPR & AUROC & AP \\
\midrule
Prompt LR & 22.9 $\pm$ 0.0 & 3.3 $\pm$ 0.0 & 74.4 $\pm$ 0.0 & 80.1 $\pm$ 0.0 \\
Prefix TF--IDF LR & 28.2 $\pm$ 0.0 & 1.6 $\pm$ 0.0 & 81.1 $\pm$ 0.0 & 86.3 $\pm$ 0.0 \\
Current-state LR & 43.7 $\pm$ 0.0 & 5.6 $\pm$ 0.0 & 83.1 $\pm$ 0.0 & 86.5 $\pm$ 0.0 \\
Running-mean LR & 46.1 $\pm$ 0.0 & 6.2 $\pm$ 0.0 & 83.5 $\pm$ 0.0 & 86.5 $\pm$ 0.0 \\
Dynamic GRU & 38.1 $\pm$ 6.3 & 3.3 $\pm$ 1.0 & 83.6 $\pm$ 0.9 & 87.6 $\pm$ 0.8 \\
\bottomrule\end{tabular}
\end{table}


\paragraph{检测器权衡。}表~\ref{tab:main}完整呈现基线比较。静态探针、MLP和循环头在共同校准目标下形成互补的召回率—误报率运行特征，为应用选择检测器提供依据。

\begin{table}[t]
\centering\footnotesize
\caption{域内告警位置；位置统计以正确检出为条件，召回率分母保留所有不安全回答。}\label{tab:alert_positions}
\begin{tabular}{lrrrr}\toprule
Method & Recall & Mean $\tau$ & Med. $\tau$ & $\tau/T$ (\%) \\
\midrule
Prompt LR & 22.9 $\pm$ 0.0 & 8.0 $\pm$ 0.0 & 8.0 $\pm$ 0.0 & 14.8 $\pm$ 0.0 \\
Prefix TF--IDF LR & 28.3 $\pm$ 0.0 & 20.7 $\pm$ 0.0 & 8.0 $\pm$ 0.0 & 26.2 $\pm$ 0.0 \\
Current-state LR & 43.9 $\pm$ 0.0 & 27.5 $\pm$ 0.0 & 16.0 $\pm$ 0.0 & 34.6 $\pm$ 0.0 \\
Running-mean LR & 46.3 $\pm$ 0.0 & 15.1 $\pm$ 0.0 & 8.0 $\pm$ 0.0 & 23.5 $\pm$ 0.0 \\
Dynamic GRU & 38.4 $\pm$ 6.1 & 47.9 $\pm$ 15.1 & 41.8 $\pm$ 18.3 & 54.0 $\pm$ 11.5 \\
\bottomrule\end{tabular}
\end{table}


\begin{table}[t]
\centering\footnotesize
\caption{提示分组配对自助法：Recall@128差值及95\%区间（百分点）。}\label{tab:uncertainty}
\begin{tabular}{lp{.34\columnwidth}rr}\toprule
Split & Contrast & $\Delta$ & 95\% CI \\
\midrule
test & Dynamic GRU vs. Prompt LR & +15.2 & [+12.7, +17.8] \\
test & Dynamic GRU vs. Current-state LR & -5.6 & [-7.1, -3.9] \\
test & Dynamic GRU vs. Prefix TF--IDF LR & +9.9 & [+7.6, +12.2] \\
external & Dynamic GRU vs. Prompt LR & +12.3 & [+6.7, +18.8] \\
external & Dynamic GRU vs. Current-state LR & -11.7 & [-20.6, -3.6] \\
external & Dynamic GRU vs. Prefix TF--IDF LR & +12.3 & [+6.7, +18.8] \\
\bottomrule\end{tabular}
\end{table}


\paragraph{同题不同结局。}测试集中同一问题同时存在两种最终标签的子集包含97组、207条回答。冻结阈值下，动态GRU与当前状态逻辑回归的Recall@128分别为20.1\%和22.3\%，误报率分别为6.1\%和6.7\%。这一同题子集刻画了两种检测器在共享问题、不同回答结局下的表现。

\paragraph{历史聚合对照。}历史均值逻辑回归的Recall@128为46.1\%，动态GRU相对差值为-8.1个百分点。该对照量化因果历史积累的作用，补充当前状态与学习式循环模型的比较。

\paragraph{不确定性。}表~\ref{tab:uncertainty}保留提示分组配对自助法区间。每次配对抽样在检测器之间保留一致的提示分组和已拟合种子。

\paragraph{预设对照。}动态GRU减去静态探针的Recall@128差值为-5.64个百分点，提示分组95\%区间为[-7.10, -3.93]。静态探针侧重召回覆盖，GRU则具有较低实测误报率，两类运行特征对应不同监测需求。
